\documentclass[11pt,a4paper]{article}
\usepackage{fontspec,amsmath,amssymb,amsthm,microtype}
\setmainfont{lmroman10-regular.otf}[BoldFont=lmroman10-bold.otf,ItalicFont=lmroman10-italic.otf,BoldItalicFont=lmroman10-bolditalic.otf]
\setmonofont{lmmono10-regular.otf}[Scale=MatchLowercase]
\usepackage[margin=25mm]{geometry}
\usepackage[unicode,hidelinks]{hyperref}
\setlength{\emergencystretch}{3em}
\setlength{\parskip}{0.4em}
\setlength{\parindent}{0pt}
\allowdisplaybreaks[2]
\raggedbottom
\title{\texorpdfstring{The actual KLMC2 Wasserstein theorem and its source corrections}{The actual KLMC2 Wasserstein theorem and its source corrections}}
\author{LangevinLean mathematical exposition}
\date{}
\begin{document}\maketitle
\noindent Source module: \texttt{Kinetic\allowbreak{}Langevin/\allowbreak{}Canonical\allowbreak{}KLMC2.\allowbreak{}lean}. The accompanying HTML pairs this mathematical exposition with the exact Lean statements and records the scope of the formalization.


This is the complete second-order sampling endpoint for the actual implemented four-kernel Gaussian chain. The continuous reference, its true Gibbs law and Gaussian moments, fresh-noise independence, local Taylor and memory error, stochastic Hessian comparison, and output-law identities are proved internally. Taking the true initial-coupling infimum yields the desired initial Wasserstein term.


\section{\texorpdfstring{Take the true transport infimum with the exact dimension term}{Take the true transport infimum with the exact dimension term}}


For every finite-cost initial plan, the checked canonical chain theorem bounds its actual L\ensuremath{{}^{2}} position discrepancy from a stationary reference. The output law is independent of that plan. The reference endpoint has the actual Gibbs law. The L\ensuremath{{}^{2}} transport-cost inequality therefore gives the same bound for the actual output Wasserstein distance. A finite-output theorem excludes an infinite distance before passing to real values. The all-couplings affine-infimum lemma then replaces each plan’s cost by the true initial \(W_2\). With \(p=|\iota|\), \(q=1-mh/(4\gamma)\), \(C_p=M_2\sqrt{p^2+2p}+M\sqrt{Mp}\), and \(\alpha=(m/(4M_2h))^2\), the result is \(W_2(\nu_n,\pi)\le\sqrt2q^nW_2(\nu_0,\pi)+\sqrt2h^2C_p/m+8(M-m)h e^{-(\alpha-p)/8}/m\). The exact Gaussian fourth moment remains valid for every finite coordinate set, including dimension one.



{\raggedright\noindent Lean declarations: \texttt{klmc2\_\allowbreak{}wasserstein\_\allowbreak{}bound\_\allowbreak{}exact}.\par}

\section{\texorpdfstring{Retain the source footnote’s sharper tail}{Retain the source footnote’s sharper tail}}


The source dimension restriction implies \(\alpha\ge5p\). Hence \((\alpha-p)/8\ge\alpha/10=m^2/(160M_2^2h^2)\). Multiply this exponential comparison by the nonnegative actual tail coefficient and apply the exact sampler bound. This gives the genuine chain theorem with the sharper inverse-step-square Gaussian tail, not merely a scalar observation detached from the algorithm.



{\raggedright\noindent Lean declarations: \texttt{klmc2\_\allowbreak{}wasserstein\_\allowbreak{}bound\_\allowbreak{}sharp\_\allowbreak{}tail}.\par}

\section{\texorpdfstring{Recover the printed bound in its valid dimension range}{Recover the printed bound in its valid dimension range}}


For \(p\ge2\), \(\sqrt2\sqrt{p^2+2p}\le2p\), and \(\sqrt2\sqrt{Mp}=\sqrt{2Mp}\). The local bias becomes \(2M_2ph^2/m+Mh^2\sqrt{2Mp}/m\). Also \(M-m\le M\) and \(e^{-(\alpha-p)/8}\le e^{-p/2}\). These comparisons produce the supplied Theorem 3 formula. Its proof’s hidden dimension restriction is stated explicitly. In dimension one the preceding exact-moment theorem supplies the valid full result. The underlying comparison includes both stochastic Hessian-difference terms omitted by the supplied displayed subtraction; those terms were repaired through the stationary free-OU bundle.



{\raggedright\noindent Lean declarations: \texttt{klmc2\_\allowbreak{}wasserstein\_\allowbreak{}bound}.\par}

\end{document}
